\documentclass{situationist-preprint}

\title{Choosing What to Want}
\msid{ACT-14-a}
\msfamily{ACT-14}
\msdiscipline{Action, Agency and Free Will}
\mstype{Principal manuscript}
\mskeywords{desire; second-order desire; identification; voluntary control; habituation; attention; self-formation}

\begin{document}
\maketitle

\begin{abstract}
We hold people answerable for their desires, praise those who have learned to want the right things, and speak of choosing to want a different life. Yet no one can start wanting a saucer of mud by deciding to. We argue that we can choose our desires in part. Second-order desires, desires about which first-order desires to have and to be moved by, can be formed by choice; first-order desires cannot, but they follow second-order choices at a distance, through the management of attention, habit and circumstance. The first claim rests on the nature of desire: a first-order desire is a state that responds to features of its object that make it seem good, and a state produced directly by a decision, in response to reasons for having it, is not of this kind. The second rests on the object of a second-order desire, which is the agent's own motivation, so that the question it answers is a practical question about what to make of oneself, which decision can settle. The third is a datum about the plasticity of desire. Together they give an indirect, partial and lagged control over what we want, parallel to the control over belief that attention gives.
\end{abstract}

\section{The problem}

\citet{Anscombe1957} observed that one cannot simply want a saucer of mud: to want something one must see it under some aspect in which it is desirable, and someone who says he wants a saucer of mud, with nothing further to say about why, has not described a want we can understand. One cannot, it seems, adopt a desire by fiat. Yet ordinary moral thought treats desires as partly up to us. Aristotle held that we become temperate by doing temperate acts, so that the person who acts from a vicious appetite is responsible for having become the kind of person who has it \citep[II.1, III.5]{AristotleNE}. \citet{Frankfurt1971} distinguished first-order desires, desires to do or have things, from second-order desires, desires to have or not to have certain first-order desires, and the subset of these he called second-order volitions, desires that a given first-order desire be one's will. A person, on his view, is a being capable of second-order volitions. \citet{Taylor1976} argued that we are responsible for ourselves through the strong evaluations by which we articulate what we care about, and \citet{Callard2018} has argued that much valuable agency consists in aspiration, the process of coming to want what one does not yet fully want.

These views suggest that we have some control over our desires, but they leave open how much and of what kind. We argue that the control is direct at the second order and indirect at the first, and that the first-order desires follow the second-order choices at a distance. The structure is parallel to the one EPI-11 defends for belief: belief is not under direct voluntary control, attention is, and belief follows attention.

\section{Framework}

\begin{definition}[Orders of desire]\label{d:orders}
A \emph{first-order desire} is a desire whose object is an action or state of affairs not concerning the agent's own desires. A \emph{second-order desire} is a desire whose object is the agent's having, lacking, or being moved by a first-order desire. A second-order desire that a first-order desire be effective in action is a \emph{second-order volition} \citep{Frankfurt1971}.
\end{definition}

\begin{definition}[Direct and indirect control]\label{d:control}
An agent has \emph{direct control} over an attitude if she can form it by deciding to, in response to considerations she takes to favor having it. She has \emph{indirect control} if she can, by deciding to do other things, bring it about that she has it.
\end{definition}

\begin{principle}[Responsiveness of desire]\label{p:resp}
A state counts as a first-order desire for $x$ only if it is formed and sustained in response to features of $x$ that present it, to the agent, as in some way good or attractive.
\end{principle}

\begin{principle}[Plasticity]\label{p:plastic}
First-order desires change in response to what the agent attends to, to what she repeatedly does, and to the circumstances she places herself in; these are within her control.
\end{principle}

Principle~\ref{p:resp} is Anscombe's point, made general: desire is a way of being drawn to something under an aspect of its desirability. It is consistent with the view that desires have evaluative content, and with the view that they are dispositions to act together with the appearance of reasons, and it does not require that the aspect be correctly perceived. Principle~\ref{p:plastic} records a familiar fact. A person who gives up sugar finds, after some months, that she no longer wants sweet things as she did; a person who attends to the suffering of factory-farmed animals finds that she no longer wants to eat their meat; a person who takes up running comes to want to run. Habituation and attention are the mechanisms Aristotle identified, and modern work on self-control adds the selection of situations.

\section{The argument}

\begin{lemma}\label{l:first}
First-order desires are not under direct voluntary control.
\end{lemma}

\begin{proof}
Suppose an agent forms, by deciding, a state directed at $x$. A decision is formed in response to considerations the agent takes to favor deciding so. If she is offered a reward for wanting $x$, the reward is such a consideration, but it is a consideration in favor of \emph{having the desire}, not a feature of $x$ that makes $x$ attractive: the reward is paid for wanting the saucer of mud, and the mud remains as unattractive as before. A state formed directly by the decision is therefore formed in response to considerations about having it, not in response to features of $x$, and by Principle~\ref{p:resp} it is not a first-order desire for $x$. If instead the considerations are features of $x$ that make it attractive, then attending to them may produce the desire; but then the desire is produced by the features attended to, not by the decision, and the control is indirect.
\end{proof}

The distinction in the proof is the one \citet{Hieronymi2006} draws between evaluative control, which we exercise over an attitude by settling the question it answers, and managerial control, which we exercise by manipulating it as an object. A first-order desire answers to the question whether its object is attractive, and that question is not settled by deciding.

\begin{lemma}\label{l:second}
Second-order volitions are under direct voluntary control.
\end{lemma}

\begin{proof}
The object of a second-order volition is the agent's own motivation: that a given first-order desire be, or not be, effective in her action. The question it answers is whether to be moved by that desire, which is a practical question about what to do with one's own motivations, and the considerations bearing on it are considerations in favor of being so moved: the value of the desire's object, its fit with the agent's other ends, the kind of person she would become. A decision is formed in response to just such considerations. When the agent decides, on reflection, that her desire to smoke shall not move her, she forms a second-order volition by deciding, and the decision responds to the considerations that bear on the question the volition answers. By Definition~\ref{d:control}, she has direct control over it.
\end{proof}

A second-order volition so formed is close to an intention, and shares its commitment. ACT-06 argues that intentions are plans with commitment that do not reduce to beliefs and desires; a second-order volition is a plan for one's own motivation, a policy of treating certain desires as reasons and not others. \citet{Bratman1996} analyzes identification with a desire in these terms, as a decision to treat it as providing a justifying reason in deliberation. The proof does not extend to every second-order desire. A person may find, without deciding anything, that she wishes she did not care so much about others' opinions; such a wish is a second-order desire that is found rather than chosen. What is chosen is the volition that settles what to do about it.

\begin{lemma}\label{l:follow}
Second-order volitions bring about corresponding first-order desires indirectly, partially and with a delay.
\end{lemma}

\begin{proof}
A second-order volition that a desire for $x$ be effective, or that a desire for $y$ cease to be, commits the agent to policies: to attend to the features of $x$ that make it attractive and away from those of $y$, to act repeatedly in the ways the desire for $x$ would prompt, to avoid situations that excite the desire for $y$. These are actions within her control, which she can decide to perform. By Principle~\ref{p:plastic}, first-order desires change in response to them. The change is indirect, since it runs through attention and habituation (Definition~\ref{d:control}). It is partial, since Principle~\ref{p:plastic} says that desires change in response to these policies, not that they change as far as the agent wills, and some desires, such as those rooted in addiction or temperament, resist reshaping. It is delayed, since habituation takes time and attention works by accumulation.
\end{proof}

\begin{theorem}\label{t:main}
We can choose our desires in part: we choose our second-order desires, in the form of second-order volitions, and first-order desires follow at a distance.
\end{theorem}

\begin{proof}
By Lemma~\ref{l:second}, second-order volitions are under direct control, so they can be chosen. By Lemma~\ref{l:first}, first-order desires are not. By Lemma~\ref{l:follow}, first-order desires change in response to the policies that second-order volitions commit us to, indirectly, partially and with delay, that is, at a distance.
\end{proof}

The parallel with EPI-11 is exact in structure. There, belief is not voluntary because a belief must respond to considerations bearing on its truth, while attention is voluntary, and belief follows attention. Here, first-order desire is not voluntary because it must respond to the attractions of its object, while second-order volition is voluntary, and desire follows the policies of attention and habituation that volition commits us to. In both cases control is exercised through the agent's direction of her own mind, and in both it grounds a responsibility that is indirect: responsibility for belief through responsibility for inquiry, responsibility for desire through responsibility for what one has made of oneself. ACT-08 argues that such self-formation need not start from nothing; Theorem~\ref{t:main} describes one way it proceeds.

\section{Objections and replies}

\begin{objection}[Regress]
If a second-order volition is authoritative because the agent chooses it, the choice must express a further desire, of the third order, and so on. Nothing at any order has the authority Frankfurt's account requires \citep{Watson1975}.
\end{objection}

\begin{reply}
The objection assumes that a choice is the expression of a desire of the next order up. On Lemma~\ref{l:second}, the volition is formed by a decision responsive to reasons for being so moved, and its authority derives from those reasons, not from a further desire. This is Watson's own point, made against Frankfurt: what gives some desires authority is the agent's evaluative standpoint, her judgment of what is worth pursuing \citep{Watson1975}. We accept it. The second-order volition is the agent's taking up that standpoint into a commitment about her own motivation, and no regress begins.
\end{reply}

\begin{objection}[The unwilling addict]
Frankfurt's unwilling addict has a second-order volition against his craving, and the craving moves him anyway. Choosing second-order desires gives no control worth the name.
\end{objection}

\begin{reply}
The addict shows that control is partial, as Theorem~\ref{t:main} says. The thesis is not that first-order desires follow second-order volitions reliably, but that they follow them at a distance. Even for the addict, the volition commits him to policies, treatment, avoidance of cues, the cultivation of other attachments, that change his desires over time in many cases. Where they do not, the limit is set by Principle~\ref{p:plastic}, and the agent is in the position of one who acts against her better judgment, as ACT-04-a describes, or one who is compelled.
\end{reply}

\begin{objection}[Inauthenticity]
A desire one has cultivated deliberately is not really one's own; authentic desires are those one finds in oneself.
\end{objection}

\begin{reply}
The objection would make the temperate person's moderate appetites less her own than the glutton's. Desires shaped by habituation are the agent's in the way a skill acquired by practice is hers. \citet{Callard2018} argues that coming to value something one does not yet fully value is a central form of rational agency, and the person who succeeds has the resulting desires as fully as anyone has desires. Being one's own, in the sense of ACT-01-a, depends on the mechanism being one's own and responsive to reasons, not on the absence of a history of self-shaping.
\end{reply}

\begin{objection}[Second-order desires are found]
We discover what we want to want, as we discover what we want. Second-order desires are no more chosen than first-order ones.
\end{objection}

\begin{reply}
We granted that some second-order desires are found. But the volition that settles whether a desire shall move one is not merely found, since it answers a practical question about what to do, and practical questions are settled by deciding. One may discover that one wishes to be less ambitious, and then decide whether to act on the wish; the decision is the volition, and it is chosen.
\end{reply}

\section{Conclusion}

We cannot want something by deciding to, since a first-order desire answers to the attractions of its object and not to reasons for having it. We can decide which of our desires shall move us, since that is a practical question about our own motivation. Such decisions commit us to policies of attention, practice and circumstance, through which our first-order desires change, slowly and incompletely. We choose what we want to want, and what we want follows at a distance.

\begin{thebibliography}{99}\small
\bibitem[Anscombe(1957)]{Anscombe1957} Anscombe, G.~E.~M. (1957). \emph{Intention}. Basil Blackwell.
\bibitem[Aristotle(NE)]{AristotleNE} Aristotle. \emph{Nicomachean Ethics}. Trans. T. Irwin, 2nd ed. Hackett, 1999.
\bibitem[Bratman(1996)]{Bratman1996} Bratman, M.~E. (1996). Identification, decision, and treating as a reason. \emph{Philosophical Topics}, 24(2), 1--18.
\bibitem[Callard(2018)]{Callard2018} Callard, A. (2018). \emph{Aspiration: The Agency of Becoming}. Oxford University Press.
\bibitem[Frankfurt(1971)]{Frankfurt1971} Frankfurt, H.~G. (1971). Freedom of the will and the concept of a person. \emph{The Journal of Philosophy}, 68(1), 5--20.
\bibitem[Hieronymi(2006)]{Hieronymi2006} Hieronymi, P. (2006). Controlling attitudes. \emph{Pacific Philosophical Quarterly}, 87(1), 45--74.
\bibitem[Taylor(1976)]{Taylor1976} Taylor, C. (1976). Responsibility for self. In A.~O. Rorty (ed.), \emph{The Identities of Persons}. University of California Press.
\bibitem[Watson(1975)]{Watson1975} Watson, G. (1975). Free agency. \emph{The Journal of Philosophy}, 72(8), 205--220.
\end{thebibliography}

\end{document}
